\section{Nesting and the Joint-Completion Distinction}

\paragraph{Proof of Proposition~\ref{prop:nesting}.}
Take $F=(t,\Pi)\in\F(\T,C,q)$. Since $t_i\in\T_i\subseteq\T'_i$, the
same completion is allowed after coarsening. Every event keeps its
positive-overlap graph and its occupancy is at most $C\le C'$. Core coverage,
buffer nonreuse, and $q$ remain unchanged. Hence
$F\in\F(\T',C',q)$. Set inclusion gives nesting of reachable counts and their
maxima whenever defined. Minimizing a fixed functional over a superset weakly
lowers its minimum; maximizing weakly raises its maximum.

\paragraph{Proof of Proposition~\ref{prop:jointcompletion}.}
There is no $AC$ overlap. Thus the three-vertex overlap graph can be connected
only if both $AB$ and $BC$ are edges under the same $s$.
Given $s\in[0,3]$, positive $AB$ overlap requires $s<1$, whereas positive $BC$
overlap requires $s+1>3$, or $s>2$. The requirements are inconsistent.
Individually they are feasible, for example at $s=0$ and $s=3$, respectively.
The obstruction is joint temporal feasibility, not simultaneous capacity.

\paragraph{Outer envelopes can err in either direction.}
Replacing each uncertain interval by its union over supports creates
$I_B^{\mathrm{out}}=[0,4)$ in the preceding example. Together with $A$ and $C$
this is connected and has maximum occupancy two. The envelope model therefore
admits an event with no valid common completion.

For the reverse direction, take three length-one intervals
$I_i=[s_i,s_i+1)$ with independent supports $s_i\in[0,2]$.
The choices $(s_1,s_2,s_3)=(0,3/4,3/2)$ form a connected chain at capacity
two. All three outer envelopes equal $[0,3)$, whose simultaneous occupancy is
three; the capacity-two envelope model rejects this feasible event. Hence
full-envelope substitution is neither an inner nor an outer event model in
general. These examples identify the necessary constraint distinction; they
do not measure how often it changes an aggregate in a public dataset.

\paragraph{Capacity-equivalent information loss.}
For fine support $\T$ and coarser $\T'\supseteq\T$, an auxiliary descriptive
quantity is
\begin{equation}
 \kappa_{\T\to\T'}(C)
 =\min\{k\in\mathbb Z_{\ge0}:M(\T,C+k)\ge M(\T',C)\},
 \label{eq:kappa}
\end{equation}
with value $+\infty$ if no such capacity exists. This definition presumes the
displayed maxima are defined. It compares reachable support, not outcome
identification. Replacing the coarse maximum by a certified feasible lower
bound $L\le M(\T',C)$ can only lower the required $k$, provided the fine-time
maxima used for the comparison are themselves certified. Such a comparison
is a lower bound on $\kappa$, not its sharp value.

\section{Two-Member Matching and Sensitivity Proofs}

\paragraph{Proof of Theorem~\ref{thm:k2matching}.}
Map every two-member event partition to its set of edges. Disjoint membership
makes this a matching; compulsory core assignment saturates every core.
Conversely, such a matching defines disjoint admissible pairs, with unmatched
buffers unused. These maps are inverse on partitions. When temporal
admissibility is defined by independent row supports, disjoint pairs use
disjoint timestamp variables, so their individual witnesses combine into one
completion. This assertion would require further verification in the presence
of cross-event constraints. It is not a bijection of timestamp-completed
worlds, since a given partition can admit many completions.

The selected rows number $|C_0|+q=2m$. Thus fixing buffer support $q$ is
equivalent to fixing matching size $m$, when this size is integral and
feasible. Pair additivity gives equality of objective values term by term.
Any declared normalization constant must be applied to both endpoints.

\paragraph{Proof of Theorem~\ref{thm:k2alternating}.}
Each vertex in $M^-\mathbin{\triangle}M^+$ has degree at most two and at most
one incident edge of each matching. Every nontrivial component is consequently
an alternating path or even cycle. Cycles have equal color counts; paths have
imbalance $-1$, $0$, or $1$. Because $|M^-|=|M^+|=m$, path imbalances sum to
zero. Common edges cancel from $w(M^+)-w(M^-)$; grouping the remaining terms by
component and dividing by $m$ proves Eq.~\eqref{eq:alternating-width}.
For the core-cover feasible sets, a path endpoint cannot be a core: every core
has degree one in both matchings. Balanced component exchanges preserve core
coverage and cardinality, while imbalanced paths must be exchanged in
cardinality-balanced collections. Additional side constraints, when present,
may prevent exchanging components separately. Different optimal witnesses can
give different decompositions but have the same endpoint values.

\paragraph{Proof of Proposition~\ref{prop:k2omission}.}
Only the right-hand side increases, which proves nesting at fixed $q$.
Every edge outside $E_0$ has positive core incidence, so $\Gamma=0$ forces all
selected edges into $E_0$. In a core-saturating matching,
$\sum_{e\in M}|e\cap C_0|=|C_0|$. The budget is therefore redundant at
$\Gamma=|C_0|$. Optimization over nonempty nested feasible sets gives the
endpoint inequalities.

\paragraph{Proof of Proposition~\ref{prop:k2thresholdcontraction}.}
Every feasible matching on $E'$ is feasible on $E$, giving the interval
inclusion. The decision rule is $\ind\{H\ge\eta\}$, so a nonempty feasible set
is ambiguous precisely when it contains an attained value below $\eta$ and
one at least $\eta$. This is equivalent to the stated endpoint condition.
A narrower real-valued range may still straddle the same threshold.

\section{Fixed-Span Connectivity and Total Unimodularity}
Let exact rational intervals have endpoints in the sorted finite list
$a=u_0<u_1<\cdots<u_\ell=b$ inside a proposed span. Intervals not contained
in $[a,b)$ are excluded. For the remaining candidates define
\[
 A_{ki}=\ind\{s_i\le u_k,\ e_i\ge u_{k+1}\},\qquad
 B_{ki}=\ind\{s_i<u_k<e_i\}.
\]
Here $k=0,\ldots,\ell-1$ indexes segment rows $A$, and
$k=1,\ldots,\ell-1$ indexes boundary rows $B$. Fix a core root $r$ and distinct
companion $c$. For additive weights $w_i$ the fixed-span LP is
\begin{align*}
 \max_x\quad &\sum_i w_ix_i\\
 \text{s.t.}\quad &1\le (Ax)_k\le C
       && k=0,\ldots,\ell-1,\\
 &(Bx)_k\ge1 &&k=1,\ldots,\ell-1,\\
 &0\le x_i\le1,\qquad x_r=x_c=1.
\end{align*}
Replacing max by min changes none of the arguments.

For an integral selection, the segment lower bounds imply that its union
fills the proposed span. A connected positive-overlap family must cover these
segments and must strictly cross every internal endpoint boundary: without a
crossing interval, rows to the left and right cannot be joined by a
positive-overlap path. Conversely, an interval crossing a boundary overlaps
the intervals covering the adjacent elementary segments. Inducting from left
to right puts all segment-covering intervals in one component. Each selected
interval has positive length and overlaps that component on a segment it
covers. Thus the constraints characterize connectivity. Segment upper bounds
characterize occupancy, including half-open boundary times because their
active set equals that of the segment immediately to the right.

Interleave the rows as $A_0,B_1,A_1,\ldots,B_{\ell-1},A_{\ell-1}$.
An interval beginning at $u_p$ and ending at $u_q$ occupies the consecutive
block $A_p,B_{p+1},\ldots,A_{q-1}$. This binary matrix has the consecutive-ones
property by columns and is totally unimodular. Signed row duplication for
lower and upper inequalities preserves total unimodularity, as does appending
identity rows for bounds and fixed membership. Integer right-hand sides then
make every vertex integral. This proves Theorem~\ref{thm:oracle}; arbitrary
pair or global cardinality constraints are not part of this TU statement.

Every feasible event has a span starting and ending at observed endpoints,
so $O(|V|^2)$ span choices suffice. Requiring one companion is an exhaustive
way to enforce at least two members, using at most $|V|-1$ LPs per root and
span. With rational weights, each LP has polynomial encoding size. For a span
strictly larger than the root interval, coverage already forces a companion;
the implementation exploits this to skip unnecessary cases. Multiple roots
may generate the same membership column, which is deduplicated.

\section{Full-Master and Branch-and-Price Certificates}

\paragraph{Proof of Proposition~\ref{prop:cg}.}
The restricted-master dual satisfies the constraints for generated columns.
If exact pricing finds no positive reduced gain, those same dual variables
satisfy every omitted-column inequality as well. They are full-master dual
feasible. The restricted primal solution is also full-master feasible, and
the two have equal objectives by restricted-master strong duality. Weak
duality therefore proves full-master optimality. A restricted-master value
alone does not provide this certificate.

\paragraph{Proof of Proposition~\ref{prop:bp}.}
At each closed branch node, exact compatible pricing establishes the full
node-LP optimum. If a free optional row has fractional usage $z_i\in(0,1)$,
every integer solution satisfies either $z_i=0$ or $z_i=1$; these are
exhaustive children. Once all row usages are integral, a fractional
distinct-column solution admits active rows $i,j$ with
\[
 0<\sum_{R\ni i,j}\lambda_R<1.
\]
To see why, suppose instead that all pair usages were zero or one. For an
active row $i$ of total usage one, all positive columns containing it would
have identical co-membership with every other row and hence identical member
sets. After deduplication this gives one column of weight one, contradicting
fractionality.

An integer partition places $i,j$ together or separately. For a priced column,
together permits both-in or both-out; separate forbids both-in.
Enumerating the consistent Boolean cases fixes individual variables without
adding arbitrary rows to the TU formulation. This yields exhaustive exact
pricing at each child, although the number of cases need not be polynomial.

The node relaxation upper-bounds all integer solutions below that node.
An incumbent is a lower bound after replay of core coverage, buffer nonreuse,
connectivity, and capacity. With distinct columns and finite disjunctions,
the full search is finite. If every branch is fathomed, the incumbent is
optimal, or infeasibility is certified if none exists. At early termination
all live regions, including the currently processed node, retain valid upper
bounds; unfinished pricing uses a previously established ancestor bound or a
trivial bound such as $|B_0|$. Taking the maximum of these bounds and the
incumbent gives the global upper bound. Discarding an unpriced live region,
or treating its restricted LP value as an upper bound, would invalidate it.

\section{Master Nonintegrality Witness}
At capacity two, use four cores
$c_0=[1,2)$, $c_1=[4,7)$, $c_2=[0,2)$, and $c_3=[1,2)$, and five buffers
$b_0=[0,7)$, $b_1=[5,6)$, $b_2=[5,7)$, $b_3=[1,2)$, and $b_4=[2,6)$.
Complete event-column enumeration has LP value four. One fractional solution
uses weights $1/2,1,1/2,1/2,1/2,1/2$ on
\[
\{c_0,c_2\},\ \{c_1,b_1\},\ \{c_2,b_3\},\ \{c_3,b_3\},\
\{c_0,b_0,b_4\},\ \{c_3,b_0,b_4\}.
\]
Every core has total coverage one and each buffer has use at most one. Complete
integer enumeration gives value three, attained for example by
$\{c_0,c_2\}$, $\{c_3,b_0,b_1\}$, and $\{c_1,b_4\}$.
The branch-and-price implementation closes this witness in 17 processed nodes,
using both buffer-usage and pair-co-membership branching.

\section{Controlled-Truth Generator and Evaluation Contract}
Each controlled instance has three core rows and eight candidate buffers. Core
intervals are staggered. Every core receives a direct true member, and one or
two cores receive a sequential member whose interval overlaps the direct
member. Distractors are temporally plausible and have broad public outcome
values. The generated truth has peak occupancy at most two; the same instances
can therefore be evaluated under $C=2,3,4$, where higher $C$ enlarges only the
analyst's feasible set.

The composition benchmark reveals the true selected-buffer count $q$ but never
the selected identities. This isolates relation ambiguity from support-count
error. The feasible point comparator maximizes an outcome-blind temporal score
over all worlds with that $q$. A nearest-$q$ rule ignores feasibility.
Maximum-cardinality, maximum-overlap pair matching and positive-overlap
connected components are support baselines that do not receive $q$.

For threshold $\eta$, the frontier certifies a positive decision when
$\underline H\ge\eta$, a negative decision when $\overline H<\eta$, and
ambiguity otherwise. If a feasible point world disagrees with the generated
truth while the frontier is identified, two feasible worlds would lie on
opposite sides despite both endpoints lying on one side, a contradiction.
This explains the observed 100\% ambiguity flag rate for feasible-point errors;
the experiment measures how often such errors occur.

Candidate truncation ranks rows without using the public outcome. ``True world
representable'' requires every true member to remain and the exact true
partition to exist in the truncated master. ``Frontier exists at true $q$''
only requires some alternative world with the same count. ``Scalar covered''
is weaker still: the true mean can fall between endpoints by coincidence even
when the true joint world is absent. These nested validation concepts are
reported separately.

\section{Annotated-Relation Audit Detail}

\paragraph{Annotation and replay contract.}
We retain timestamp, planar position, speed, and direction of motion. Labels
determine eligibility and quarantine positive IDs in partial, nonpositive,
multiple-ID, one-sided, or conflicting annotations rather than silently treat
them as singletons. Candidate edge scoring remains label-free. Before an
endpoint is published, replay checks the raw solution's integrality, degree,
cardinality, objective, and primal--dual residuals. The 195 cells with at most
12 vertices also agree with an independent subset-recursion oracle.

\paragraph{Removing oracle cardinality.}
Here $m$ denotes the annotated number of dyads, distinct from the selected
buffer count $q$ in the core--buffer model. We union fixed-$m$ frontiers over
either $m\pm1$ or every positive cardinality
through half the visible population. The latter removes labels from endpoint
cardinality, although labels still determine eligibility, exclusions, and
evaluation. Relative to annotated $m$, the all-positive-$m$ regime makes 826
rather than 806 cells feasible because some candidate graphs admit a smaller
matching than the annotated count. This is apparent feasibility, not restored
relational coverage. Mean width increases from $1.176$ to $1.297$ meters, and
certified threshold decisions fall from $1286/2418$ (53.2\%) to $1136/2478$
(45.8\%). All 20 false certificates occur where candidate support omits the
annotated world; none occurs under representable truth.

\paragraph{Cross-day, multi-query point reconstruction.}
We train a class-balanced logistic dyad score on DIAMOR-1 using only distance,
heading difference, and speed difference, freeze it, and construct
cardinality-$m$ matchings on DIAMOR-2. The same matching is evaluated for mean
distance, speed gap, and heading gap. Across 381 exact cells with representable
truth, the learned and closest-distance rules have similar mean dyad recall
(89.24\% versus 89.62\%) but different downstream errors. The learned rule
makes 44, 70, and 35 threshold errors for the three queries, compared with 32,
44, and 38 for closest-distance; every error is frontier-ambiguous. The linear
scorer is a transparent cross-day baseline, not a state-of-the-art trajectory
model.

\paragraph{Alternating-witness detail.}
In the frozen follow-up support, the 73 multi-component witness differences
contain 243 alternating paths and three cycles, with at most four edges per
component. The $k=2$ degree cap contracts 21 upper endpoints by 0.177 meters on
average among changed cells (maximum 0.577 meters) while preserving every lower
endpoint and all 246 threshold classifications. These counts illustrate
Theorem~\ref{thm:k2alternating}; they are not required for it.

\section{Frozen NYC Decision-Panel Detail}
The panel predeclares 20 eight-core windows and four 16-core stress windows.
Seventeen eight-core and all four 16-core windows are eligible. The three
outcome-blind ineligible windows are July weekend PM, October weekend AM,
and October weekend PM, all at eight cores; each failed
to produce an integrity- and cap-qualified core. All 24 terminal reports are
observed and technical failures are zero.

\begin{table*}[t]
\caption{Frozen NYC panel by core size, outcome, and capacity. Ambiguity is at
the candidate-median threshold; point split is disagreement among four
feasible deterministic point methods.}
\label{tab:paneldetail}
\small
\centering
\begin{tabular}{@{}llrrrrrr@{}}
\toprule
Core & Outcome & $C$ & Cells & Closed & Ambig. & Point split & Median width \\
\midrule
8 & Miles   & 2 & 17 & 13 & 17 & 2 & 19.291 \\
8 & Miles   & 3 & 17 & 16 & 17 & 2 & 20.441 \\
8 & Miles   & 4 & 17 & 17 & 17 & 2 & 20.058 \\
8 & Minutes & 2 & 17 & 14 & 17 & 6 & 60.879 \\
8 & Minutes & 3 & 17 & 14 & 17 & 6 & 62.657 \\
8 & Minutes & 4 & 17 & 17 & 17 & 6 & 62.881 \\
16 & Miles   & 2 & 4 & 2 & 3 & 0 & 19.626 \\
16 & Miles   & 3 & 4 & 0 & 4 & 0 & -- \\
16 & Miles   & 4 & 4 & 2 & 4 & 0 & 17.889 \\
16 & Minutes & 2 & 4 & 2 & 4 & 0 & 68.202 \\
16 & Minutes & 3 & 4 & 2 & 4 & 0 & 68.321 \\
16 & Minutes & 4 & 4 & 2 & 4 & 0 & 68.321 \\
\bottomrule
\end{tabular}
\end{table*}

A cell is certified ambiguous when attained inner feasible values lie on both
sides of the threshold; exact endpoints are not required for that conclusion.
The sole unresolved median-threshold decision is the 16-core miles cell at
$C=2$. Candidate-quartile and transparent 5-mile/30-minute references are
reported in the artifact and are not claimed operational policy thresholds.
Table~\ref{tab:paneldetail} separates numerical endpoint closure from witnessed
ambiguity. Table~\ref{tab:capacitydetail} retains the earlier common-support
comparison: the certified $C=2$ maximum of 72 selected buffers is held fixed
across capacities on one eight-core input.

\begin{table*}[t]
\caption{Frontiers of mean selected-buffer miles and trip minutes at a common
support of 72 selected buffers.}
\label{tab:capacitydetail}
\centering
\begin{tabular}{@{}crrrr@{}}
\toprule
$C$ & Miles range & Width & Minutes range & Width \\
\midrule
2 & 3.440--9.679 & 6.239 & 17.779--37.147 & 19.368 \\
3 & 3.131--15.448 & 12.317 & 15.983--55.712 & 39.729 \\
4 & 3.116--16.166 & 13.050 & 15.924--58.906 & 42.982 \\
\bottomrule
\end{tabular}
\end{table*}

\section{Frozen Branch-and-Price Scale Detail}
Table~\ref{tab:scaledetail} gives the final support-maximization lattice.
Its integer closures refer to the declared numerical and witness-replay
protocol; they do not certify the decision panel's separate outcome endpoints.
\begin{table*}[t]
\caption{All 18 predeclared NYC integer scale cells. A bracket is the valid
[integer incumbent, global upper bound] interval at termination.}
\label{tab:scaledetail}
\small
\centering
\begin{tabular}{@{}rrrrlrrr@{}}
\toprule
Core & Buffer & $C$ & Root LP & Status & Bracket & Nodes & Seconds \\
\midrule
4  & 12 & 2 & 4  & Exact & [4,4]   & 1 & 0.48 \\
4  & 12 & 3 & 8  & Exact & [8,8]   & 4 & 6.98 \\
4  & 12 & 4 & 12 & Exact & [12,12] & 2 & 6.22 \\
6  & 18 & 2 & 8  & Exact & [8,8]   & 1 & 4.04 \\
6  & 18 & 3 & 14 & Exact & [14,14] & 8 & 56.32 \\
6  & 18 & 4 & 18 & Exact & [18,18] & 1 & 15.26 \\
8  & 24 & 2 & 12 & Exact & [12,12] & 1 & 12.21 \\
8  & 24 & 3 & 20 & Exact & [20,20] & 3 & 62.05 \\
8  & 24 & 4 & 24 & Exact & [24,24] & 4 & 79.94 \\
10 & 30 & 2 & 16 & Exact & [16,16] & 1 & 52.47 \\
10 & 30 & 3 & 26 & Exact & [26,26] & 1 & 86.36 \\
10 & 30 & 4 & 30 & Exact & [30,30] & 8 & 1031.00 \\
12 & 36 & 2 & 20 & Exact & [20,20] & 1 & 145.94 \\
12 & 36 & 3 & 32 & Exact & [32,32] & 1 & 198.91 \\
12 & 36 & 4 & 36 & Exact & [36,36] & 11 & 2075.49 \\
16 & 48 & 2 & 28 & Exact & [28,28] & 1 & 415.59 \\
16 & 48 & 3 & 44 & Exact & [44,44] & 3 & 1068.73 \\
16 & 48 & 4 & 48 & Exact & [48,48] & 7 & 4193.20 \\
\bottomrule
\end{tabular}
\end{table*}

\section{Chicago Transport and Graph-Cover Pin}
The successful live audit is workflow run \texttt{33837186969}, source commit
\nolinkurl{b2e549e7e4cc674a7a880dc7789ee5f3c960d2b0}. The ZIP artifact SHA-256 is
\nolinkurl{98b7117b14c3150df655d88f171be6f05d0774449af18d354fac98f639e1226b};
the JSON report SHA-256 is
\nolinkurl{26b682d967716cfc356cfc4d68a39446bdc85c941d676bb7bea5323e1f71dca3}.
The extraction strategy is
\nolinkurl{NARROW_OVERLAP_INDEX_THEN_EXACT_START_AND_ENDPOINT_BIN_SHARDS}.
Duplicate public identifiers with inconsistent payloads fail closed, and each
shard is counted before and after retrieval. The graph certificate remains
conditional on outer released-time envelopes and is not a hidden-run closure
certificate.

\section{Artifact Claim Boundary}
The repository contains deterministic fixed-time oracles, numerical
existential-time programs, branch-and-price validators, controlled-truth
benchmarks, public-data snapshot/count contracts, and aggregate evidence
manifests. Supported claims are limited to declared candidate universes and
models. No artifact output should be interpreted as recovered partner identity,
a complete private event, an estimate of realized capacity, or a statement
about a city's proprietary production system.
